\newglossaryentry{agente}{
    name={Agente},
    description={Entidad computacional que observa el estado de su entorno y ejecuta acciones según una política determinada con el objetivo de maximizar su recompensa a largo plazo}
}

\newglossaryentry{entorno}{
    name={Entorno},
    description={El sistema dinámico con el que interactúa el agente. Responde a las acciones del agente produciendo transiciones de estado y emitiendo señales de recompensa}
}

\newglossaryentry{estado}{
    name={Estado},
    description={Representación cuantitativa o cualitativa de la situación actual del entorno en un instante de tiempo específico. Sirve como base para que el agente seleccione su próxima acción}
}

\newglossaryentry{observacion}{
    name={Observación},
    description={Subconjunto de información del estado completo del entorno que es efectivamente visible o medible por un agente individual, especialmente relevante en escenarios de observabilidad parcial}
}

\newglossaryentry{estado_global}{
    name={$s$ (estado global)},
    description={Representación conjunta de las variables relevantes del entorno en un instante determinado. En el modelo multi-agente puede incluir la información de todas las intersecciones y sus corrientes vehiculares}
}

\newglossaryentry{observacion_local}{
    name={$o_i$ (observación local)},
    description={Información disponible para el agente $i$ durante la ejecución descentralizada. Se obtiene a partir de los sensores o variables observables de la intersección que controla}
}

\newglossaryentry{accion}{
    name={$a_i$ (acción)},
    description={Decisión ejecutada por el agente $i$ en un paso temporal. En el control semafórico representa la fase o la modificación de fase seleccionada por el controlador}
}

\newglossaryentry{recompensa_t}{
    name={$r_t$ (recompensa)},
    description={Señal escalar recibida en el paso temporal $t$ después de ejecutar una acción. Se utiliza para orientar la actualización de la política o de la función de valor}
}

\newglossaryentry{gamma}{
    name={$\gamma$ (factor de descuento)},
    description={Parámetro en el intervalo $[0,1)$ que pondera la contribución de las recompensas futuras al retorno esperado}
}

\newglossaryentry{q_aprendizaje}{
    name={$Q(s,a)$ (valor de acción)},
    description={Retorno esperado de ejecutar la acción $a$ en el estado $s$ y continuar posteriormente de acuerdo con una política determinada}
}

\newglossaryentry{v_aprendizaje}{
    name={$V(s)$ (valor de estado)},
    description={Retorno esperado al comenzar en el estado $s$ y seguir una política determinada, promediando las acciones que dicha política selecciona}
}

\newglossaryentry{politica_parametrizada}{
    name={$\pi_\theta$ (política parametrizada)},
    description={Distribución de probabilidad sobre las acciones condicionada por el estado u observación y controlada por el vector de parámetros entrenables $\theta$}
}

\newglossaryentry{ventaja_estimada}{
    name={$\hat{A}_t$ (ventaja estimada)},
    description={Estimación de cuánto mejor o peor fue la acción observada en el paso $t$ respecto del valor esperado bajo la política vigente}
}

\newglossaryentry{recompensa}{
    name={Recompensa},
    description={Señal escalar devuelta por el entorno que evalúa cuán buena o mala fue la acción ejecutada por el agente en un estado particular. El objetivo del agente es maximizar la suma total de estas señales}
}

\newglossaryentry{politica}{
    name={Política (Aprendizaje por Refuerzo)},
    description={Estrategia que define el comportamiento del agente en un momento dado, mapeando los estados observados a probabilidades de selección de las acciones disponibles}
}

\newglossaryentry{actor}{
    name={Actor (Aprendizaje por Refuerzo)},
    description={Componente de la arquitectura del modelo de aprendizaje responsable de seleccionar y ejecutar la acción en el entorno basándose en el estado actual y la política que está aprendiendo}
}

\newglossaryentry{critico}{
    name={Crítico (Aprendizaje por Refuerzo)},
    description={Componente encargado de estimar la función de valor, prediciendo la recompensa futura esperada para evaluar la calidad de las acciones tomadas por el Actor y así guiar su actualización}
}

\newglossaryentry{funcion_valor}{
    name={Función de Valor},
    description={Estimación matemática de la cantidad total de recompensa que un agente puede esperar acumular en el futuro, comenzando desde un estado particular y siguiendo una política específica}
}

\newglossaryentry{ventaja}{
    name={Ventaja (Aprendizaje por Refuerzo)},
    description={Métrica que cuantifica cuánto mejor o peor resulta tomar una acción particular en un estado dado, en comparación con el desempeño promedio esperado según la política actual}
}

\newglossaryentry{mdp}{
    name={Proceso de Decisión de Markov},
    description={Marco matemático utilizado para modelar la toma de decisiones en situaciones donde los resultados son en parte aleatorios y en parte bajo el control de un agente tomador de decisiones}
}

\newglossaryentry{fase}{
    name={Fase (Ingeniería de Tránsito)},
    description={Estado temporal de un controlador semafórico que habilita el derecho de paso exclusivo a una o más corrientes vehiculares incompatibles para que crucen una intersección de forma segura}
}

\newglossaryentry{ciclo}{
    name={Ciclo (Ingeniería de Tránsito)},
    description={Secuencia completa y repetitiva de todas las fases semafóricas programadas en una intersección particular}
}

\newglossaryentry{flujo}{
    name={Flujo vehicular},
    description={Cantidad de vehículos que pasan por una sección transversal específica de un carril o calzada durante un intervalo de tiempo determinado}
}

\newglossaryentry{densidad}{
    name={Densidad vehicular},
    description={Número de vehículos presentes en un tramo o longitud específica de un carril en un instante dado}
}

\newglossaryentry{cola}{
    name={Cola vehicular},
    description={Fila de vehículos que se detienen o avanzan muy lentamente debido a una restricción en la capacidad de la vía, tal como la luz roja de un semáforo}
}
